\documentclass[UTF8,a4paper,11pt]{ctexart}
\usepackage[margin=2.5cm]{geometry}
\usepackage[numbers,sort&compress]{natbib}
\usepackage[hidelinks]{hyperref}
\title{矿山遥感变化检测的泛化与可靠性：研究证据综述}
\author{}
\date{2026年9月5日}
\begin{document}
\maketitle
\begin{abstract}
矿山遥感变化检测正由单一区域的变化提取转向跨地区可复用的方法评估。公开矿山足迹及双时相变化数据为这一转变提供了基础，但数据规模、检测精度与预测可靠性分别回答不同问题。现有变化表征和域泛化方法已经能够抑制部分外观干扰，概率校准与选择性预测也提供了处理不确定输出的工具；这些进展仍需在标签边界清楚、矿区独立、复核预算一致的条件下组合验证。矿山场景的研究增量应落实到背景误报与细小变化保留之间的关系，以及可靠性排序对剩余错误的影响。以真实误差诊断约束方法选择，并区分离线复核模拟和真实人员试验，是形成可信结论的关键。
\end{abstract}

\section{引言}
矿山足迹反映采掘、堆放及相关设施在地表形成的空间范围，变化图则关注两个时相之间的差异。全球矿山用地资料拓展了可追溯的空间参照，但静态边界不提供变化发生时间和全部变化成因\cite{Maus2020Mining}。更新版全球矿山用地产品进一步扩展了空间覆盖，其不同版本之间的差别仍需结合影像时间、解释规则和制图过程理解\cite{Maus2022Mining}。因此，应用研究首先需要明确标签所代表的观测对象，再讨论网络结构和泛化能力。

矿山活动导致的可见变化可以为环境监测提供线索，但单凭二值影像掩码难以区分污染、沉降、生产安全风险及其因果关系。方法评价应围绕变化位置是否准确、错误在何处集中、有限复核能够纠正多少错误展开。这一定位允许将变化检测与人工判读衔接，也限定了算法结论的使用范围。

\section{公开矿山数据与标签边界}
MineNetCD 为全球矿山双时相变化检测提供了统一任务和模型比较基础\cite{Yu2024MineNetCD}。其原文明确采用矿区层面的训练、验证与测试划分，因而“矿区隔离”已是基准的一部分。后续研究需要核对下载版本、原始场景与裁切块的对应关系，并说明新增协议与原协议的差别。单纯重新生成划分清单不能构成方法创新。

EuroMineNet 将欧洲矿山足迹研究扩展至多年 Sentinel-2 观测\cite{Yu2026EuroMineNet}。年度足迹标签可支持相应时相的制图和按定义构造的变化任务，但与高分辨率双时相变化数据之间存在输入、时间间隔和标签语义差异。比较两个数据集时，应先在各自任务中建立基线；如果开展迁移，必须交代重采样、波段选择、标签映射和目标域数据使用方式。不同任务的 IoU 差值不是单一域偏移效应。

全球铜矿资料提供矿山边界和属性信息，可辅助选择有色金属矿山案例\cite{Li2025Copper}。其中静态用地类别能够帮助解释某一年份的矿山空间组成，却不能直接替代其他年份的变化真值。矿种专题分析应以名称、坐标及年份匹配为前提，并报告匹配失败和覆盖不足，不把少量便利案例扩展成全部有色矿区的代表性证据。

\section{变化表征的进展与背景约束}
孪生全卷积网络提供了双时相特征提取与差异融合的清晰基线，适合检验数据配对、标签读取和训练流程\cite{Daudt2018Siamese}。SNUNet-CD 通过密集连接组织多层特征，为定位信息保留提供了另一类结构参照\cite{Fang2022SNUNet}。这类方法有助于识别网络容量与特征融合的作用，但仅与一个较弱模型比较，不能证明新方法解决了跨矿区泛化问题。

双时相图像变换器 BIT 通过紧凑语义表征建模时空上下文\cite{Chen2022BIT}，ChangeFormer 将层级变换器编码器与多层感知机解码器结合，体现了另一种长距离关系建模方式\cite{Bandara2022ChangeFormer}。ChangeMamba 则将状态空间模型用于多时相交互\cite{Chen2024ChangeMamba}。这些工作表明可选骨干已经丰富；新课题应把骨干作为受控实现条件，防止容量、预训练和训练预算的差别被误认为核心机制贡献。

FCCDN 已在特征提取和融合过程中引入约束\cite{Chen2022FCCDN}。因此，泛泛提出“增加一致性损失”不足以形成清晰增量。矿山变化任务值得追问的是，约束应该施加在什么区域，以及约束减少背景误报时是否同时抹去了细小变化。若使用训练标签确认的未变化区域，应进一步排除标签边界和无效像素，比较全图约束、背景约束、随机等面积约束及普通监督训练。在相同训练预算下，只有局部误差结构随约束对象发生可重复的变化，才能支持对机制的解释。

\section{域泛化与可比实验}
域泛化研究需要在目标数据未参与训练和模型选择时评价迁移能力。DomainBed 的研究提醒，经验风险最小化在控制实验条件和模型选择后可以成为有竞争力的对照\cite{Gulrajani2021DomainBed}。WILDS 则从不同真实任务展示了训练环境与测试环境变化带来的评价困难\cite{Koh2021WILDS}。两者并不直接证明某个矿山数据集必然存在相同程度的失效，却为设置独立测试和公平对照提供了依据。

MixStyle 通过训练样本的特征统计混合增加风格多样性，是成本较低的域泛化参照\cite{Zhou2021MixStyle}。DonaNet 面向遥感伪变化，引入域相关风格移除和跨时相一致性学习\cite{Zang2025DonaNet}。CDDGNet 进一步研究变化特征与域特征的解耦\cite{Zhao2026CDDG}。这些直接近邻说明，本课题不能把“未见区域检测”或“变化不变性”称为尚无人研究的问题。可检验的差别应是矿山标签语义下的约束区域、背景混淆分层和对细小变化的代价。

AWDA 使用无标签目标域数据开展跨数据集变化检测适应\cite{AWDA2025}。这一设定与完全不访问目标域数据的泛化不同，结果必须分栏报告。官方预训练权重也需要按训练范围核查：若权重可能已见过自建测试矿区，则可用于复现原模型，但不能据此声称该测试矿区对模型完全未知。

\section{概率校准、错误排序与复核效用}
温度缩放为预测概率提供了简单的后处理校准参照\cite{Guo2017Calibration}。它不等于错误识别器；概率刻度改变后，样本排序可能基本不变。已有研究说明，数据分布变化时，训练分布上良好的校准未必保持\cite{Ovadia2019Uncertainty}。同时，后续研究也发现较新模型未必比早期模型校准更差\cite{Minderer2021Calibration}。这些结果要求将“高置信错误多发”写成待验证假说，而不是所有网络都满足的先验事实。

建筑变化检测已经研究标签和模型层面的不确定性\cite{Li2025Uncertainty}。矿山研究可以借鉴其问题意识，但建筑目标、矿山裸露地物和年度足迹的语义并不一致。对本课题而言，Brier 分数、变化类分层误差、可靠性曲线与区域错误排序应共同报告；整体校准指标可能受到大量背景像素影响，不能单独代表变化类可靠性。

SelectiveNet 联合学习预测与拒绝，明确呈现风险和覆盖的权衡\cite{Geifman2019SelectiveNet}；按类别构造选择性预测的方法进一步说明，各类错误的控制方式可以不同\cite{Gangrade2021Selective}。拒绝选项的理论研究则区分了代价、风险与覆盖目标\cite{Franc2023Reject}。这些方法学结果不能直接转译成矿山场景的分布无关保证，但提示复核实验必须固定预算和分母。

实际复核单元应覆盖整幅有效影像，包括预测为未变化的区域，否则无法发现高置信漏检。以不重叠空间窗口为例，可在固定面积比例下比较随机选择、原始不确定度、温度缩放和候选错误排序分数。除了保留区域的平均错误率，还要报告全图尺度的剩余误报、剩余漏检、送审面积以及真实变化被送审的比例。利用既有标签模拟“送审区域全部纠正”只能估计理想复核收益，真实复核者仍可能漏判，实际工作时间也未必与面积成正比。

\section{讨论：空间相关与统计推断}
空间相关数据的交叉验证需要与实际预测目标一致，随机样本划分和空间外推衡量的是不同能力\cite{Roberts2017CV}。生态制图研究给出了忽略空间结构会过度乐观的实例\cite{Ploton2020Spatial}；该现象不能作为 MineNetCD 已存在泄漏的证据。blockCV 为构造空间或环境分离折提供了工具思路\cite{Valavi2019BlockCV}，最近邻距离匹配验证进一步强调验证样本与实际预测空间的距离关系\cite{Mila2022Nearest}。矿区研究应首先明确是新矿区预测、相邻区域插值还是不同观测条件迁移，再选择合适的验证办法。

大量切片和像素并不等于同等数量的独立场景。重复训练种子主要反映优化随机性，矿区间差异则反映场景变化；把二者混成所有像素的标准误会高估精度。可先按种子汇总每个矿区指标，再按矿区进行成对重采样，同时独立报告种子离散度。极少变化矿区、空变化矿区和数据缺失矿区需要预先规定分母和缺失值处理方式。

\section{展望}
最有价值的下一步是用真实误差限制研究假说。先复现原协议，审计矿区及权重来源，在独立开发数据上判断背景误报、边界偏差和细小变化漏检是否普遍出现。随后只围绕一个主要失效比较局部背景约束与直接近邻方法，并把检测改善和复核收益分开检验。若来源可追溯性不足，结论应收缩到已验证场景；若方法改善只能由容量或重复调参解释，则应放弃相应机制主张。

扩展数据和有色矿案例应作为适用边界检验，而非预先承诺全面迁移。对于规模有限的开放课题，公开数据审计、一个有明确对照的核心改进及一套复核评价足以构成连贯工作，前提是结果可复现且失败条件被如实报告。

\section{结论}
现有证据支持将矿山变化检测组织为数据语义、变化判别和复核效用三个相互衔接的层次。矿区隔离、特征一致性、域泛化和概率校准均已有研究基础，本课题的贡献需要通过具体约束对象和错误代价的比较来确立。当前最稳妥的研究路径是保护已有基准的可比性，在独立矿区上检验背景约束能否保留真实变化，并验证可靠性排序是否在同一复核预算下减少剩余错误。

\bibliographystyle{gbt7714-nsfc}
\bibliography{mining_change_references}
\end{document}
